% SPDX-License-Identifier: GPL-3.0-or-later
% Copyright (C) 2026 Alessandro Greco
\chapter{匿名、监控与元数据}\label{cap:01-anonimato}
\aggiornatoal{2026 年 9 月 23 日}{Tor Metrics 数据截至 2026 年 9 月 20 日}

\begin{obiettivi}
学完本章后，你将能够：
\begin{itemize}[nosep]
  \item 区分机密性、匿名性、假名性、不可链接性和不可观察性；
  \item 解释为什么通信元数据能够揭示与内容同样多、甚至更多的信息；
  \item 描述流量分析，并说明为什么仅靠加密无法阻止它；
  \item 说明哪些人在使用 Tor、为什么使用，以及用户群体的多样性为何本身就是一种保护。
\end{itemize}
\end{obiettivi}

本书讨论的是一种旨在保护特定信息的系统：它主要保护的不是我们在网上\emph{说了什么}，而是我们\emph{在与谁通信}。在研究系统的工作方式之前，必须先弄清楚它试图解决的问题，以及这个问题为什么不同于我们更熟悉的“加密”问题。

% ---------------------------------------------------------------------------
\section{内容与元数据}\label{sec:anon-metadati}

当两台计算机通过互联网通信时，每个数据包都会携带两类信息。\textbf{内容}是应用程序实际交换的数据，例如消息文字或网页内容。\textbf{元数据}描述的则是通信\emph{本身}：谁在发送、发给谁、什么时候发送、发送得多频繁，以及数据量有多大。有些元数据是网络正常运行所必需的，因为路由器必须知道应把数据包送到哪里。目的 IP 地址以及用于接收回复的源 IP 地址，都会以明文形式出现在每个数据包中。

如今 HTTPS 已十分普及，加密因此能够保护内容，但元数据仍然可见。任何能够观察网络的人——从用户的互联网服务提供商，到大型互联网交换中心的运营者——仍可能看到用户连接了哪个网站、连接发生在什么时候、持续了多久以及产生了多少流量。

\begin{figure}[htbp]
\centering
% SPDX-License-Identifier: GPL-3.0-or-later
% Copyright (C) 2026 Alessandro Greco
%% 简体中文版 EPUB 替代文本。
\altfig{HTTPS 连接示意图：通信内容位于标注为已加密的方框中，而周围的元数据仍以明文可见，包括源 IP 地址、目的 IP 地址、时间信息以及数据包大小和发送频率。}
% HTTPS 能保护什么，以及观察者仍能看到什么。
\begin{tikzpicture}[font=\sffamily\scriptsize,
  riga/.style={draw,minimum height=0.65cm,align=center},
  visibile/.style={riga,fill=avviso!12,draw=avviso!60},
  protetto/.style={riga,fill=torverde!15,draw=torverde!60!black}]
  \node[client,minimum width=1.6cm] (u) at (0,0) {用户};
  \node[server,minimum width=1.6cm] (s) at (10.6,0) {网站};
  \draw[very thick,torgrigio!50] (u) -- (s);
  \node[font=\sffamily\scriptsize,text=torgrigio] at (5.3,0.45) {HTTPS};
  \node[visibile,text width=1.9cm] (a) at (2.5,-1.5) {发送方\\IP 地址};
  \node[visibile,text width=1.9cm] (b) at (4.6,-1.5) {网站\\IP 地址};
  \node[visibile,text width=1.9cm] (c) at (6.7,-1.5) {时间与持续时长};
  \node[visibile,text width=1.9cm] (d) at (8.8,-1.5) {发送与接收\\的字节数};
  \node[protetto,text width=4.2cm] (e) at (3.4,-2.9) {请求的页面与内容};
  \node[protetto,text width=4.2cm] (f) at (7.9,-2.9) {凭据、消息正文};
  \node[text=avviso] at (0.5,-1.5) {\textbf{元数据}};
  \node[text=avviso,font=\sffamily\tiny] at (0.5,-1.85) {明文};
  \node[text=torverde!45!black] at (0.5,-2.9) {\textbf{内容}};
  \node[text=torverde!45!black,font=\sffamily\tiny] at (0.5,-3.25) {已加密};
  \foreach \n in {a,b,c,d} \draw[dotted,avviso!70] (\n.north) -- (\n.north |- u.south);
\end{tikzpicture}

\caption{观察者在 HTTPS 连接上能够看到什么。内容受到加密保护，但元数据仍然可见；Tor 试图保护的正是这类信息。}
\label{fig:metadati}
\end{figure}

\subsection{元数据会暴露什么}

人们很容易把元数据看成无关紧要的信息：\enquote{他们只知道我和谁说话，却不知道我说了什么。}这种直觉并不可靠。几个例子就能说明原因：
\begin{itemize}
  \item 某人在凌晨两点访问家庭暴力援助机构的网站，随后又访问离婚律师的网站；
  \item 在一篇调查报道发布前几天，一名记者持续与该调查涉及的某政府部门员工交换消息；
  \item 一名员工经常访问工会网站，或者某种特定疾病患者协会的网站。
\end{itemize}
在这些情形中，无须读取通信内容，就已经能够推断出大量信息。对收集者而言，元数据还有另一个优势：结构化程度高，易于大规模存储和分析，也便于与其他数据交叉关联。曾先后担任 NSA 和 CIA 局长的 Michael Hayden 在 2014 年 4 月 1 日约翰斯·霍普金斯大学的一场讨论中，用一句后来广为流传的话概括了这一点：\enquote{我们会根据元数据杀人}（本书译文），随后立刻补充说\enquote{但我们不会用这些元数据做那件事}~\cite{hayden2014metadata}。后一句限定并没有改变核心事实：收集元数据的人认为，这些信息本身就足以支撑极其重大的判断。

% ---------------------------------------------------------------------------
\section{流量分析}\label{sec:anon-traffico}

\textbf{流量分析}通过研究通信元数据来推断信息，而不必读取通信内容。它的历史几乎与电信本身一样悠久。即使在二十世纪的战争中，仅仅观察敌方无线电台多久发一次报、与谁通信、通信节奏如何，也可能在完全无法解密消息的情况下推断出一支军队的组织结构。

在互联网上，流量分析主要有两种形式：
\begin{description}
  \item[地址分析] 链路上的观察者可以看到源 IP 地址和目的 IP 地址。这是最简单的一类，也是 Tor 最直接处理的问题：网络中不存在一个天然的单点能够同时看到两个端点。
  \item[流量大小与时序分析] 即使看不到完整的地址信息，观察者也可以按数据量和时间模式比较不同流。如果某用户在 10:15:02 下载了 3\,MB，而网络另一处的一台服务器也在 10:15:02 发出了 3\,MB 数据，那么两次观察很可能属于同一通信。这种\textbf{流量关联}是 Tor 的主要局限，本书会反复回到这个问题（\cref{cap:04-modello-minaccia,cap:26-attacchi}）。
\end{description}

% ---------------------------------------------------------------------------
\section{精确的术语}\label{sec:anon-termini}

“匿名”一词经常被用于不同含义，而这些含义之间的混淆，在匿名系统的发展史上曾导致不少设计错误。科学文献，尤其是 Pfitzmann 和 Hansen 提出的术语体系~\cite{pfitzmann2010terminology}，会区分以下几个概念：
\begin{description}
  \item[机密性] 未授权的一方无法读取内容。这是加密主要提供的属性。
  \item[匿名性] 在一组主体中无法识别某个主体。这组主体称为\textbf{匿名集}。匿名从来不是绝对属性，它总是“相对于其他人”存在。匿名集越大、成员越多样，保护通常越强。
  \item[假名性] 一个人使用并非真实姓名的标识符活动，而且这个标识符可能长期保持稳定。使用虚构名称的账户是一个假名，而不等于匿名：该账户的所有行为仍可能被彼此关联。
  \item[不可链接性] 观察者无法判断两个关注对象——例如两条消息、一个用户与一个网站，或两次访问——是否相关。Tor 试图在用户与其目的地之间提供的，正是这种属性。
  \item[不可观察性] 观察者甚至无法判断通信\emph{是否正在发生}。这是强得多的属性，而 Tor 本身并不提供它：监视用户网络连接的人通常能够看到该用户正在使用 Tor，除非同时使用可插拔传输（\cref{cap:16-censura}）。
\end{description}

\begin{storia}[把身份识别与路由分离]
1995 年发明洋葱路由的三位 Naval Research Laboratory 研究人员——David Goldschlag、Michael Reed 和 Paul Syverson——更愿意避免把目标简单描述成“匿名”。他们在 1996 年的一篇论文中写道，他们的目标是\enquote{把身份识别与路由分离}；Syverson 后来回顾那段历史时再次强调了这一点~\cite{syverson2011peel}。例如，一个用户可能希望向正在通信的系统（如办公室计算机）证明自己的身份，但并不希望网络中的每一个中间节点都知道“位于这个地点的人正在连接那台办公室计算机”。
\end{storia}

% ---------------------------------------------------------------------------
\section{谁需要匿名？}\label{sec:anon-utenti}

根据 Tor Metrics 的估算，2026 年 8 月和 9 月期间，每天有约 280 万至 350 万用户在不使用网桥的情况下连接 Tor 网络，平均约 320 万。意大利的估算值每天约为 \num{44000} 至 \num{83000}，平均约 \num{60000}。这些数字存在相当大的不确定性，因为网络只能估算用户数量，而不是对用户逐一计数（\cref{cap:24-metrics}）。这些用户是谁？他们来自非常不同的群体：
\begin{itemize}
  \item 不希望广告商、互联网服务提供商或所使用网络的运营者追踪自己的普通用户；
  \item 记者及其消息来源、活动人士和人权捍卫者，特别是生活在威权政府统治国家的人；
  \item 身处网络审查环境、借助 Tor 访问被屏蔽网站的人（\cref{cap:16-censura}）；
  \item 需要隐藏自身位置的暴力或迫害受害者；
  \item 不希望向竞争对手暴露研究兴趣的公司；
  \item 需要在访问网站或收集信息时不暴露身份的执法和军事人员——洋葱路由最初正是在美国海军实验室中为这类需求而诞生。
\end{itemize}
正如任何通信工具一样，也有一小部分人会利用 Tor 从事违法活动。公共讨论往往把注意力集中在这一群体上。第四部分以及 \cref{cap:25-legale} 会基于证据和具体案例讨论这个问题。

\subsection{匿名需要同伴}

用户多样性的重要性还有更深一层原因。如果一种匿名系统只由某一个群体使用——例如只有军方使用——那么它反而无法很好地保护该群体：观察者只要看到某人连接这张网络，就会立即知道这个人属于军方；看到从网络中出来的流量，也会知道它来自军方。Syverson 在讨论美国海军原本可以建设的一张专用网络时，明确指出了这一问题：所有从该网络流出的通信都会被识别为海军流量~\cite{syverson2011peel}。每个人的匿名性都取决于还有多少其他用户，以及这些用户彼此之间有多么不同。概括这一原则的一句话——同时也是 Dingledine 和 Mathewson 一篇论文的标题——是\enquote{匿名需要同伴（anonymity loves company）}~\cite{dingledine2006company}。这也是 Tor 被设计成向所有人开放、免费且开源的原因之一：使用者越多，它越有可能更好地保护每一个人。

% ---------------------------------------------------------------------------
\section*{小结}
\begin{itemize}
  \item 加密保护内容；元数据（谁、和谁、什么时候、传输了多少）仍然可见，而且可能揭示大量信息。
  \item 流量分析根据通信的地址、数据量和时序推断信息。
  \item 匿名性意味着在一个集合中无法识别某个主体；必须把它与假名性、不可链接性和不可观察性区分开。
  \item Tor 服务于数百万背景各异的用户，而这种多样性本身就是其保护机制的一部分。
\end{itemize}

% ---------------------------------------------------------------------------
\section*{练习}

\begin{esercizio}[1]\label{es:01-metadati}
列出互联网服务提供商在其客户阅读 HTTPS 新闻网站文章时可以观察到的元数据。如果客户改用 VPN，哪些信息仍然可见？如果使用 Tor，又会怎样？
\end{esercizio}

\begin{esercizio}[1]\label{es:01-termini}
将下列情形分别归类为匿名、假名或两者都不是：使用虚构名称的社交媒体账户；留在公共公告板上的未署名消息；只用姓名首字母签署的信件；投入募捐箱的现金捐款。
\end{esercizio}

\begin{esercizio}[2]\label{es:01-insieme}
在一小时内，某网站收到 50 名不同 Tor 用户的访问。其中一人发布了一条评论，包含意大利语使用者常见的拼写错误，而且发布时间是意大利当地凌晨 3 点。这会怎样改变该用户的匿名集？这个例子说明用户行为会如何影响匿名性？
\end{esercizio}

\begin{esercizio}[2]\label{es:01-vpn}
比较商业 VPN 与 Tor：在两种系统中，谁有能力把用户与其访问目的地关联起来？信任分别集中在哪里？
\end{esercizio}

% ---------------------------------------------------------------------------
\section*{延伸阅读}
\begin{itemize}
  \item A.~Pfitzmann, M.~Hansen, \titolo{A terminology for talking about privacy by data minimization}~\cite{pfitzmann2010terminology}：隐私术语的重要参考。
  \item P.~Syverson, \titolo{A Peel of Onion}~\cite{syverson2011peel}：由洋葱路由的发明者之一回顾其历史。
  \item Tor Metrics, \url{https://metrics.torproject.org/}：按国家估算的 Tor 用户数量。
\end{itemize}
